\chapter{Test e validazione}
\label{cap:test}

\section{Validazione del livello di rilevamento}
\label{sec:test-detection}

\textbf{DA AGGIORNARE}

Le scelte lasciate aperte in fase di implementazione --- se costruire il livello regex/NER sopra Microsoft Presidio (§\ref{sec:valutazione-presidio}) e quale motore di NER adottare --- sono qui chiuse su base empirica.

\paragraph{Metodo.} La valutazione riusa l'impianto di scoring del progetto: un corpus annotato di documenti sintetici in cui ogni PII è marcata con il proprio \texttt{pii\_type} e la sua posizione. Un rilevamento conta come vero positivo quando condivide il \texttt{pii\_type} con uno span di \textit{ground truth} e vi si sovrappone con \textit{Intersection over Union} $\geq 0.5$; ogni rilevamento e ogni span sono appaiati al più una volta (\textit{matching} greedy). Da qui si derivano precision, recall e F1 micro-mediati. Ogni configurazione è incapsulata dietro il contratto \texttt{PIIDetector} e misurata sullo stesso corpus, così da confrontarle sullo stesso metro.

\paragraph{Risultati.} Il confronto è nella Tabella~\ref{tab:test-detection}. \textbf{Il livello a pattern di Presidio} (espressioni regolari e checksum) supera la baseline regex proprietaria (F1 0.85 contro 0.79) con precisione perfetta e senza falsi positivi, e copre nativamente categorie che la baseline mancava; viene quindi adottato, coerentemente con la scelta di incapsularlo dietro \texttt{PIIDetector} (§\ref{sec:valutazione-presidio}). Per il \textbf{livello NER}, il modello generico spaCy si rivela inadeguato: F1 0.08, con recall quasi nullo e falsi positivi (etichetta come \textit{location} identificatori strutturati, ad esempio un IBAN). Sostituendolo con \textbf{GLiNER}, NER \textit{zero-shot} specifico per le PII, il livello NER sale a F1 0.44 e --- dato decisivo --- la detection completa (pattern + NER) passa da F1 0.77 a \textbf{0.89}: GLiNER chiude i buchi del NER, portando \texttt{address} da 0 a F1 1.00 e il recall di \texttt{person\_name} a 1.00.

\begin{table}[htbp]
	\centering
	\caption{Validazione del livello di rilevamento sul corpus annotato (P, R, F1).}
	\label{tab:test-detection}
	\small
	\begin{tabular}{|l|c|c|c|}
		\hline
		\textbf{Configurazione}                  & \textbf{P}    & \textbf{R}    & \textbf{F1}   \\
		\hline
		Baseline regex (proprietaria)            & 0.93          & 0.68          & 0.79          \\
		\hline
		Presidio --- solo pattern/checksum       & 1.00          & 0.74          & 0.85          \\
		\hline
		NER spaCy (\texttt{it\_core\_news\_lg})  & 0.17          & 0.05          & 0.08          \\
		\hline
		NER GLiNER (\texttt{gliner\_multi\_pii}) & 0.75          & 0.32          & 0.44          \\
		\hline
		Presidio (pattern + NER spaCy)           & 0.75          & 0.79          & 0.77          \\
		\hline
		\textbf{Presidio (pattern + NER GLiNER)} & \textbf{0.89} & \textbf{0.89} & \textbf{0.89} \\
		\hline
	\end{tabular}
\end{table}

\paragraph{Limiti.} Il corpus è ridotto (poche decine di span) e sintetico: i valori vanno letti come \textbf{direzionali}, non come stime definitive, e la loro conferma richiede un corpus più ampio e realistico. Restano inoltre due punti aperti: la precisione di GLiNER su \texttt{person\_name} è ancora migliorabile (over-tagging occasionale), e l'inferenza GLiNER su CPU è più lenta di spaCy --- accettabile, trattandosi di un modello \textit{encoder} e non generativo, ma da tenere presente per il livello di rilevamento ricorrente (blocco B4).
